% page 301 du fac-simile (image p-300 du scan)
% bloc 162.6 x 233.3 mm ; marge gauche 39.4 mm
\pgc{17.780mm}{0.000mm}{
\l{\cel{55}293}
\l{}
\l{\sou{konciza}. (\sou{pri stilo}).Sen to quo ne esas necesa por la senco.}
\l{\cel{3}- DEFIS.}
\l{}
\l{\sou{kondamnar}. (\sou{trans.}) I.Deklarar kulpoza, per verdikto : l.}
\l{\cel{3}frapar (ulu) per judicio qua deklaras ke lu kulpis per}
\l{\cel{3}kriminir, deliktir, detrimentir ; 2. (\sou{metaf.})(pri\cel{1}\sou{libri,}}
\l{\cel{3}\sou{letri, e}\cel{1}c.) Koaktar ad ulo puniso. - II. Deklarar ke (ulu)}
\l{\cel{3}esas kulpoza, blaminda, reprimandinda.- III.1. Deklarar}
\l{\cel{3}ke (ulo) ne plus esas uzebla. 2. Deklarar kom ne plus uzenda}
\l{\cel{3}(pordo, fenestro, quan onu klozas tale ke lu ne plus esas}
\l{\cel{3}apertebla.)- EFIS.}
\l{}
\l{\sou{kondensar}. (\sou{trans.}) Igar (gaso, vaporo) plu densa, per la plu-}
\l{\cel{3}proximigo di lua molekuli. - DEFIS.}
\l{}
\l{\sou{kondensatoro}. (\sou{elektro}).Aparato quan onu uzas por akumular,}
\l{\cel{3}sur surfaco, elektro pozitiva o negativa.}
\l{}
\l{\sou{kondicionar}. (\sou{trans., per)}\cel{1}.\cel{2}Submisar a klauzo obliganta de qua}
\l{\cel{3}dependas la valideso (di akto, e c.) - EFIS.}
\l{}
\l{\sou{kondilo}. (\sou{anat.)}\cel{2}Saliajo artikala, osta, rondigita ca-facie -}
\l{\cel{3}platigita, ta-facie. - EF.}
\l{}
\l{\sou{kondilono.}\cel{2}(\sou{patol.}) Surkreskajo morbala en la regiono di la}
\l{\cel{3}anuso o di la perineo.}
\l{}
\l{\sou{kondimento}. Substanco quan onu uzas por plufortigar la saporo}
\l{\cel{3}di nutrivi, sive lor lia manjeso an la repasto-tablo, sive}
\l{\cel{3}lor lia preparo. - DEFIS.}
\l{}
\l{\sou{kondolar}. (\sou{trans., pri, pro}).Expresar sua partopreno di trauro,}
\l{\cel{3}di desfortuno, qui afliktas (ulu). - DEFS.}
\l{}
\l{\sou{kondominiar}. (\sou{trans.}) Kunposedar o kundominacar (lando). - EF.}
\l{}
\l{\sou{kondoro}. (\sou{zool.}) Granda vulturo, en la Andi. - L.\cel{1}\sou{sarcorhamphus}}
\l{\cel{3}\sou{gryphus}. - DEFIRS.}
\l{}
\l{\sou{konduktar}. (\sou{trans.}) (\sou{fiziko}).Propagar de ca a ta punto la mani-}
\l{\cel{3}festi elektrala, kalorala. - EFIRS.}
\l{}
\l{\sou{konduktoro}. (en\cel{1}\sou{vehilo publika : tramveturo, autobuso, treno}).}
\l{\cel{3}La komizo di qua la rolo esas relatar kun la voyajeri, e c. -}
\l{\cel{3}DEFRS.}
\l{}
\l{\sou{kondutar}. (\sou{netrans.,}\cel{1}ad).Agar, en ta à ca cirkonstanco, en sua}
\l{\cel{3}vivo, ye ta o ca maniero. - EFIS.}
\l{}
\l{\sou{konektar}. (\sou{trans., kun}) Ligar strete kun kozo sama-natura. -}
\l{\cel{3}(\sou{elektro}). Inter-relatigar, per elektro-konduktili, mashini}
\l{\cel{3}od organi. - EF.}
}
